\section{Trace-Length Mechanism}
\label{sec:trace-length}
\label{sec:length}

For a closed triangle $T$ and a rectifiable target $E$, define
\[
 L_E(T):=\Haus(T\cap E).
\]
\begin{lemma}[Connected open-cover budget]
\label{lem:open-cover-budget}
Let $X$ be compact and connected, with a finite Borel measure $\mu$
positive on every nonempty relatively open subset. For a finite open cover
with at least two nonempty relative pieces,
\[
 \mu(X)<\sum_j\mu(U_j\cap X)\le\sum_j\mu(T_j\cap X).
\]
Consequently closed-piece upper bounds whose sum is at most $\mu(X)$
already contradict the open cover.
\end{lemma}
\begin{proof}
If the nonempty relative pieces were pairwise disjoint, one and the union
of the others would disconnect $X$. Some pair therefore overlaps in a
nonempty relatively open set, of positive measure. Integrating the covering
multiplicity gives strict inequality. Closure only enlarges each piece.
\end{proof}

For $X=\partial H$ or $S$, use length; for $X=H$, use area. Each measure
has the stated full-support property, and the seven distinguished roles
meet the corresponding targets in positive measure. In particular,
\begin{equation}
 \Haus(E)<L_E(T_C)+\sum_{i=0}^5L_E(T_i),\qquad E\in\{\partial H,S\},
 \label{eq:length-subadditivity}
\end{equation}
whenever these targets are covered. Empty center boundary traces are simply
omitted from the collection. The local geometric caps remain separate from
this common final budget argument.

\begin{figure}[H]
\centering
\begin{tikzpicture}[scale=.84]
  \foreach \x/\title in {0/$\partial H$: length $6$,4/$D$: length $6$,8/$S=\partial H\cup D$: length $12$}{
    \begin{scope}[xshift=\x cm]
      \coordinate (a0) at (1,0);
      \coordinate (a1) at (.5,.866);
      \coordinate (a2) at (-.5,.866);
      \coordinate (a3) at (-1,0);
      \coordinate (a4) at (-.5,-.866);
      \coordinate (a5) at (.5,-.866);
      \coordinate (o) at (0,0);
      \draw[hex] (a0)--(a1)--(a2)--(a3)--(a4)--(a5)--cycle;
      \node[panellabel,align=center] at (0,-1.38) {\title};
    \end{scope}
  }
  \draw[demand] (1,0)--(.5,.866)--(-.5,.866)--(-1,0)--(-.5,-.866)--(.5,-.866)--cycle;
  \foreach \p in {(4,0),(4.5,.866),(3.5,.866),(3,0),(3.5,-.866),(4.5,-.866)}{
    \draw[radialdemand] (4,0)--\p;
  }
  \draw[demand] (9,0)--(8.5,.866)--(7.5,.866)--(7,0)--(7.5,-.866)--(8.5,-.866)--cycle;
  \foreach \p in {(9,0),(8.5,.866),(7.5,.866),(7,0),(7.5,-.866),(8.5,-.866)}{
    \draw[radialdemand] (8,0)--\p;
  }
  \node[figlabel,align=center,text=DemandRed] at (8,-1.93)
    {the full-skeleton estimate is used only under its stated branch hypotheses};
\end{tikzpicture}

\caption{The perimeter and full-skeleton targets.  The skeleton is used only
under the hypotheses of the corresponding trace budget.}
\label{fig:strategy1-targets}
\end{figure}

\subsection{Skeleton interfaces}

\begin{lemma}[Center skeleton cap]
\label{lem:center-skeleton-cap}
If $N_E(T_C)\in\{1,2\}$ and
\[
 \left|T_C\cap\{M_0,\ldots,M_5\}\right|=1,
\]
then
\[
 L_S(T_C)<\frac32.
\]
\end{lemma}

\begin{proof}
This is the signed-center trace optimization
\zcref{lem:app-center-skeleton-cap-calculation}.
\end{proof}

\begin{lemma}[Positive-support skeleton cap]
\label{lem:positive-support-skeleton-cap}
Let $T_i$ be the closure of an original open V triangle.  If
\[
 \exists j\in\{i-1,i+1\}\quad\Haus(T_i\cap r_j)>0,
\]
then
\[
 L_S(T_i)<\frac32.
\]
\end{lemma}

\begin{proof}
Normalize to $i=0$ and translate the hexagon by $V_0$.  The five components
of $S$ on which a unit triangle containing $V_0$ can have positive trace
become the corresponding local components of the translated skeleton.  The
positive adjacent-radial trace becomes a positive boundary trace there.
Proposition~\ref{prop:unique-center-midpoint} and
Lemma~\ref{lem:center-skeleton-cap} give the bound; every nonlocal component
is more than unit distance from $V_0$, apart from zero-measure endpoints.
\end{proof}

\begin{lemma}[No-support skeleton cap]
\label{lem:no-support-skeleton-cap}
If
\[
 A_i+B_i\le1,
\]
\[
 \Haus(T_i\cap r_{i-1})=0,
 \qquad \Haus(T_i\cap r_{i+1})=0,
\]
then
\[
 L_S(T_i)\le2.
\]
\end{lemma}

\begin{proof}
See \zcref{lem:app-no-support-skeleton-calculation}; geometrically, the two
incident boundary traces and the own radial trace are the only components
that can contribute positive length.
\end{proof}

\begin{lemma}[Supercritical skeleton cap]
\label{lem:supercritical-skeleton-cap}
If $A_i+B_i>1$, then
\[
 L_S(T_i)<\frac32.
\]
\end{lemma}

\begin{proof}
The exact admissible-cell and polynomial-sign calculation is
\zcref{lem:app-supercritical-skeleton-calculation}.
\end{proof}

\begin{lemma}[Positive-support rescuer]
\label{lem:positive-support-rescuer}
Suppose $N_E(T_C)\in\{1,2\}$ and, after symmetry,
\[
 T_C\cap\{M_0,\ldots,M_5\}=\{M_0\}.
\]
If $N_+\ge2$, then for a supercritical index $s\ne0$ there is
$j\in\{s-1,s+1\}$ such that
\[
 M_s\in U_j,
 \qquad \Haus(T_j\cap r_s)>0.
\]
\end{lemma}

\begin{proof}
Choose a supercritical $s\ne0$.  Lemma~\ref{lem:self-midpoint} gives
$M_s\notin T_s$, and the displayed midpoint identity gives
$M_s\notin T_C$.  The open cover therefore puts $M_s$ in another $U_j$.
Diameter locality restricts $j$ to $s-1$ or $s+1$, and openness makes its
trace on $r_s$ have positive length.
\end{proof}

\begin{theorem}[Direct CE1/CE2 skeleton count]
\label{thm:common-skeleton-count}
If
\[
 N_E(T_C)\in\{1,2\},
\]
\[
 \left|T_C\cap\{M_0,\ldots,M_5\}\right|=1,
\]
and
\begin{equation}
 N_++N_{\mathrm{sp}}\ge3,
 \label{eq:common-skeleton-count}
\end{equation}
then the seven triangles do not cover $S$.
\end{theorem}

\begin{proof}
The three geometric caps above make the total trace strictly smaller than
$\Haus(S)=12$; the exact budget is
\zcref{lem:app-common-skeleton-count-calculation}.  This contradicts
\eqref{eq:length-subadditivity}.
\end{proof}

\subsection{Boundary trace register}

The exact evaluation is in Appendix~\ref{app:trace-length-optimization}; the
following table is the complete interface used by the routing argument.

\begin{theorem}[Boundary trace table]
\label{thm:boundary-trace-table}
For closures of original open triangles, the following bounds hold.
\begingroup
\setlength{\arraycolsep}{5pt}
\renewcommand{\arraystretch}{1.18}
\begin{equation}
\begin{array}{c|c|c}
\text{exact hypothesis}&\text{established class}&\text{exact bound}\\ \hline
N_E(T_C)=0&\CEzero&L_{\partial H}(T_C)=0\\
N_E(T_C)=1&\CEone&
 L_{\partial H}(T_C)\le\dfrac{\sqrt3}{2}-\dfrac34\\[3pt]
N_E(T_C)=2&\CEtwo&L_{\partial H}(T_C)<\dfrac12\\[3pt]
\begin{gathered}
 (o(T_i),n(T_i))\in\{(1,0),(2,0)\}\\
 A_i+B_i\le1
\end{gathered}
 &\Vdzero&L_{\partial H}(T_i)\le1\\[8pt]
\begin{gathered}
 (o(T_i),n(T_i))\in\{(1,0),(2,0)\}\\
 A_i+B_i>1
\end{gathered}
 &\Vdzero&L_{\partial H}(T_i)\le\dfrac2{\sqrt3}\\[8pt]
(o(T_i),n(T_i))\in\{(1,1),(1,2)\}
 &\Vdone\text{ or }\Vdtwo&L_{\partial H}(T_i)<\dfrac12\\
(o(T_i),n(T_i))=(2,1)&\Tlike&L_{\partial H}(T_i)<1.
\end{array}
\label{eq:boundary-trace-table}
\end{equation}
\endgroup
For every V triangle,
\begin{equation}
 L_{\partial H}(T_i)=A_i+B_i.
 \label{eq:vertex-boundary-locality}
\end{equation}
\end{theorem}

\begin{proof}
The locality identity and every row are proved in
\zcref{lem:app-boundary-trace-table-calculation}.
\end{proof}

\begin{corollary}[Branches closed by the common budgets]
\label{cor:signed-budget-branches}
The following exact alternatives are impossible.
\begin{enumerate}
\item $N_E(T_C)=0$, $N_+=1$, and $d\ge1$;
\item $N_E(T_C)\in\{1,2\}$, $N_+=0$, and $d\ge1$;
\item $N_E(T_C)=1$, $N_+=1$, and $d\ge1$;
\item $N_E(T_C)=2$, $N_+=1$, and $d\ge2$;
\item $N_E(T_C)=2$, $N_+=1$, and there is an index $i$ such that
\[
 (o(T_i),n(T_i))=(1,2)
\]
and
\[
 \{M_{i-1},M_{i+1}\}\cap T_i\ne\varnothing;
\]
\item $N_E(T_C)\in\{1,2\}$ and $N_++N_{\mathrm{sp}}\ge3$;
\item $N_E(T_C)\in\{1,2\}$ and $N_+\ge2$.
\end{enumerate}
\end{corollary}

\begin{proof}
The perimeter sums for items (1)--(5) are
\zcref{lem:app-signed-budget-calculation}.  Item (6) is
Theorem~\ref{thm:common-skeleton-count}.  For item (7), normalize the unique
C triangle midpoint and apply Lemma~\ref{lem:positive-support-rescuer}; its
rescuer is distinct from the two supercritical triangles, so item (6)
applies.
\end{proof}

\begin{proposition}[Vd2 adjacent-midpoint branch]
\label{prop:ce2-vd2-midpoint-length}
Suppose the C triangle is CE1 or CE2 and $N_+=1$. If a Vd2 role contains
an adjacent midpoint, the perimeter is not covered.
\end{proposition}
\begin{proof}
By \zcref{lem:vd2-neighbor-midpoint-cap}, that role contributes strictly
less than $1/3$. The common boundary caps give
\[
 L_{\partial H}(T_C)+\sum_iL_{\partial H}(T_i)
 <\frac12+\frac2{\sqrt3}+\frac13+4<6.
\]
Apply \zcref{lem:open-cover-budget}. The argument needs no CE2-only formula.
\end{proof}

\subsection{Routing conclusion}

The zero-gap $N_+=1$, $d\ge1$ clause is retained as an independent
alternative; the routing table now sends every zero-gap $N_+=1$ state to
\zcref{thm:reader-zero-gap-obstruction}.

\begin{proposition}[Branches closed by the trace-length mechanism]
\label{prop:length-branches}
The following routing entries and hybrid subcase are impossible.
\begin{enumerate}
 \item $N_{\mathrm{gap}}=0$ and $N_+=0$;
 \item $N_{\mathrm{gap}}=0$, $N_+=1$, and $d\ge1$;
 \item $N_E(T_C)\in\{1,2\}$ and $N_++N_{\mathrm{sp}}\ge3$;
 \item $N_E(T_C)\in\{1,2\}$, $N_+=0$, and $d\ge1$;
 \item $N_E(T_C)=1$, $N_+=1$, and $(d,t)=(1,0)$;
 \item the subcase of $N_E(T_C)=2$, $N_+=1$, and $(d,t)=(1,0)$ in which
 the unique non-$\Vdzero$ triangle $T_i$ satisfies
 \[
  (o(T_i),n(T_i))=(1,2)
 \]
 and
 \[
  \{M_{i-1},M_{i+1}\}\cap T_i\ne\varnothing.
 \]
\end{enumerate}
\end{proposition}

\begin{proof}
For item (1), Lemma~\ref{lem:gap-exhaustion} makes the incident open V
traces overlap on every edge, so
\[
 B_i+A_{i+1}>1\qquad(0\le i\le5).
\]
Summing contradicts $N_+=0$.  Items (2)--(5) are clauses (1), (6), (2), and
(3), respectively, of Corollary~\ref{cor:signed-budget-branches}; their
numerical budgets are evaluated in Appendix~\ref{app:trace-length-optimization}.
Item (6) is Proposition~\ref{prop:ce2-vd2-midpoint-length}.
\end{proof}
